\chapter{Granica funkcji}
\label{ma1:ch:function-limits}

Dotychczas badaliśmy granice ciągów, czyli odwzorowań określonych na zbiorze liczb naturalnych. Teraz rozważymy zachowanie odwzorowania, gdy jego argument zbliża się do ustalonego punktu przestrzeni metrycznej. Granica zależy od zachowania odwzorowania w pobliżu punktu, a nie od wartości w samym punkcie. Najpierw przedstawimy ujęcie metryczne, następnie szczególne rodzaje granic funkcji rzeczywistych jednej zmiennej.

\section{Punkt skupienia i definicja granicy}
\begin{definicja}
\label{ma1:def:accumulation-point}
Niech $(X,d_X)$ będzie przestrzenią metryczną, $A\subseteq X$ i $p\in X$. Punkt $p$ nazywamy punktem skupienia zbioru $A$, jeżeli dla każdego $\delta>0$ istnieje $x\in A$ taki, że $0<d_X(x,p)<\delta$. Zbiór wszystkich punktów skupienia $A$ oznaczamy przez $A'$.
\end{definicja}
Punkt skupienia nie musi należeć do $A$. Dla $A=(0,1)$ w $\R$ mamy $A'=[0,1]$. Istotnie, każdy punkt tego przedziału można przybliżać różnymi od niego punktami $A$, natomiast dla punktu spoza $[0,1]$ istnieje przedział otwarty rozłączny z $A$. Zbiór $\Z$ w $\R$ nie ma punktów skupienia: dla każdej liczby rzeczywistej można wybrać przedział, który nie zawiera żadnej liczby całkowitej różnej od jego środka.

\begin{lemat}
\label{ma1:lem:accumulation-sequence}
Niech $(X,d_X)$ będzie przestrzenią metryczną, $A\subseteq X$ i $p\in X$. Wówczas $p\in A'$ wtedy i tylko wtedy, gdy istnieje ciąg $(x_n)_{n\ge1}$ w $A\setminus\{p\}$ zbieżny do $p$.
\end{lemat}
\begin{proof}
Jeżeli $p\in A'$, to dla każdego $n\ge1$ wybieramy $x_n\in A$ z $0<d_X(x_n,p)<1/n$. Wówczas $x_n\ne p$ i $d_X(x_n,p)\to0$. Odwrotnie, jeśli $x_n\in A\setminus\{p\}$ oraz $x_n\to p$, to dla każdego $\delta>0$ od pewnego indeksu $0<d_X(x_n,p)<\delta$. Definicja \ref{ma1:def:accumulation-point} daje $p\in A'$.
\end{proof}

\begin{definicja}
\label{ma1:def:function-limit}
Niech $(X,d_X)$ i $(Y,d_Y)$ będą przestrzeniami metrycznymi, $A\subseteq X$, $f:A\to Y$ oraz $p\in A'$. Mówimy, że granicą odwzorowania $f$ w punkcie $p$ względem zbioru $A$ jest $L\in Y$, jeżeli
\[
 \forall\varepsilon>0\ \exists\delta>0\ \forall x\in A:
 \quad 0<d_X(x,p)<\delta
 \ \Longrightarrow\ d_Y(f(x),L)<\varepsilon.
\]
Piszemy $\lim_{A\ni x\to p}f(x)=L$ lub, gdy zbiór $A$ jest ustalony, $\lim_{x\to p}f(x)=L$.
\end{definicja}
Warunek $p\in A'$ zapewnia, że w każdym przekłutym otoczeniu punktu $p$ istnieją argumenty należące do $A$. Bez tego założenia implikacja z definicji byłaby dla dostatecznie małego $\delta$ spełniona dla każdego $L\in Y$. Wartość $f(p)$, o ile istnieje, nie występuje w warunku granicznym.

\begin{przyklad}
\label{ma1:ex:removable-value}
Niech $f:\R\to\R$ będzie dana wzorem $f(x)=x+1$ dla $x\ne0$ i $f(0)=7$. Dla $0<|x|<\varepsilon$ mamy $|f(x)-1|=|x|<\varepsilon$, zatem $\lim_{x\to0}f(x)=1$, mimo że $f(0)=7$. Ten sam dowód działa dla ograniczenia $f$ do $\R\setminus\{0\}$.
\end{przyklad}

\begin{twierdzenie}[Kryterium ciągowe granicy]
\label{ma1:thm:function-limit-sequential}
Niech $(X,d_X)$ i $(Y,d_Y)$ będą przestrzeniami metrycznymi, $f:A\to Y$, gdzie $A\subseteq X$, $p\in A'$ oraz $L\in Y$. Wówczas
$\lim_{A\ni x\to p}f(x)=L$ wtedy i tylko wtedy, gdy dla każdego ciągu $(x_n)_{n\ge1}$ w $A\setminus\{p\}$ zbieżnego do $p$ zachodzi $f(x_n)\to L$.
\end{twierdzenie}
\begin{proof}
Załóżmy istnienie granicy $L$ i niech $x_n\to p$, $x_n\in A\setminus\{p\}$. Dla $\varepsilon>0$ wybieramy $\delta>0$ z definicji granicy odwzorowania. Od pewnego indeksu $0<d_X(x_n,p)<\delta$, więc $d_Y(f(x_n),L)<\varepsilon$. To dowodzi zbieżności obrazów.

Odwrotnie, jeżeli warunek z definicji granicy nie zachodzi dla $L$, to istnieje $\varepsilon_0>0$ takie, że dla każdego $\delta>0$ istnieje $x\in A$ z $0<d_X(x,p)<\delta$ oraz $d_Y(f(x),L)\ge\varepsilon_0$. Dla $\delta=1/n$ wybieramy taki punkt $x_n$. Otrzymujemy ciąg zbieżny do $p$, którego obrazy pozostają w odległości co najmniej $\varepsilon_0$ od $L$. Jest to sprzeczne z założeniem.
\end{proof}
\begin{wniosek}
\label{ma1:cor:function-limit-unique}
Przy założeniach definicji \ref{ma1:def:function-limit} odwzorowanie $f$ ma w punkcie $p$ co najwyżej jedną granicę w $Y$.
\end{wniosek}
\begin{proof}
Z lematu \ref{ma1:lem:accumulation-sequence} wybieramy ciąg $x_n\in A\setminus\{p\}$ zbieżny do $p$. Gdyby $L_1,L_2$ były granicami $f$, to z kryterium ciągowego ciąg $f(x_n)$ dążyłby do obu punktów. Jednoznaczność granicy ciągu w przestrzeni metrycznej daje $L_1=L_2$.
\end{proof}
\begin{przyklad}
\label{ma1:ex:direction-dependent-limit}
Niech $A=\R^2\setminus\{(0,0)\}$ i $f(x,y)=x^2/(x^2+y^2)$. Ciągi $(1/n,0)$ i $(0,1/n)$ należą do $A$ i dążą do $(0,0)$ w metryce euklidesowej. Obrazy pierwszego ciągu są równe $1$, a drugiego $0$. Gdyby granica istniała, oba ciągi obrazów miałyby tę samą granicę. Zatem granica nie istnieje.
\end{przyklad}

\section{Ograniczenia, złożenia i kryterium Cauchy'ego}
\begin{twierdzenie}
\label{ma1:thm:limit-restriction-union}
Niech $(X,d_X)$, $(Y,d_Y)$ będą przestrzeniami metrycznymi, $f:A\to Y$, $A\subseteq X$, $p\in A'$ oraz $L\in Y$.
\begin{enumerate}[label=(\roman*)]
\item Jeśli $\lim_{A\ni x\to p}f(x)=L$ i $B\subseteq A$ oraz $p\in B'$, to $\lim_{B\ni x\to p}f(x)=L$.
\item Jeśli $A=B\cup C$, gdzie $B,C\subseteq X$ oraz $p\in B'\cap C'$, to $\lim_{A\ni x\to p}f(x)=L$ wtedy i tylko wtedy, gdy obie granice względem $B$ i $C$ istnieją i są równe $L$.
\end{enumerate}
\end{twierdzenie}
\begin{proof}
(i) Każdy punkt $B$ spełniający warunek odległości należy do $A$, więc dla tej samej pary $\varepsilon,\delta$ nierówność w przestrzeni $Y$ zachodzi na $B$.
(ii) Konieczność wynika z (i). Dla dostateczności ustalamy $\varepsilon>0$ i wybieramy $\delta_B,\delta_C>0$ dla obu granic. Jeśli $\delta=\min(\delta_B,\delta_C)$ i $x\in A$ spełnia $0<d_X(x,p)<\delta$, to $x$ należy do co najmniej jednego z tych dwóch zbiorów; w każdym przypadku $d_Y(f(x),L)<\varepsilon$.
\end{proof}
\begin{twierdzenie}
\label{ma1:thm:limit-composition}
Niech $(X,d_X),(Y,d_Y),(Z,d_Z)$ będą przestrzeniami metrycznymi, $A\subseteq X$, $B\subseteq Y$, $f:A\to B$, $g:B\to Z$, $p\in A'$, $q\in B'$ i $L\in Z$. Załóżmy, że $\lim_{A\ni x\to p}f(x)=q$, $\lim_{B\ni y\to q}g(y)=L$ oraz istnieje $\eta>0$ takie, że $f(x)\ne q$ dla wszystkich $x\in A$ z $0<d_X(x,p)<\eta$. Wówczas $\lim_{A\ni x\to p}g(f(x))=L$.
\end{twierdzenie}
\begin{proof}
Dla $\varepsilon>0$ istnieje $\rho>0$ takie, że dla $y\in B$ z $0<d_Y(y,q)<\rho$ mamy $d_Z(g(y),L)<\varepsilon$. Z granicy $f$ istnieje $\delta_0>0$ takie, że $d_Y(f(x),q)<\rho$ dla $x\in A$ z $0<d_X(x,p)<\delta_0$. Przyjmijmy $\delta=\min(\delta_0,\eta)$. Wówczas $0<d_Y(f(x),q)<\rho$ i $f(x)\in B$, więc $d_Z(g(f(x)),L)<\varepsilon$.
\end{proof}
\begin{uwaga}
\label{ma1:rem:composition-zero}
Jeżeli $q\in B$ i $g(q)=L$, założenie o unikaniu punktu $q$ można pominąć: dla $f(x)=q$ nierówność w przestrzeni $Z$ zachodzi automatycznie. Bez obu tych warunków twierdzenie jest fałszywe. Istotnie, niech $f:\R\to\R$ będzie stale równa $0$, a $g(y)=0$ dla $y\ne0$ i $g(0)=1$. Mamy $\lim_{x\to0}f(x)=0$, $\lim_{y\to0}g(y)=0$, lecz $g(f(x))=1$ dla każdego $x$.
\end{uwaga}

\begin{twierdzenie}[Kryterium Cauchy'ego granicy]
\label{ma1:thm:function-limit-cauchy}
Niech $(X,d_X)$ będzie przestrzenią metryczną, $(Y,d_Y)$ zupełną przestrzenią metryczną, $A\subseteq X$, $f:A\to Y$ i $p\in A'$. Granica $\lim_{A\ni x\to p}f(x)$ istnieje w $Y$ wtedy i tylko wtedy, gdy dla każdego $\varepsilon>0$ istnieje $\delta>0$ takie, że dla wszystkich $x,z\in A$ spełniających
$0<d_X(x,p)<\delta$ oraz $0<d_X(z,p)<\delta$ zachodzi $d_Y(f(x),f(z))<\varepsilon$.
\end{twierdzenie}
\begin{proof}
Jeżeli granica istnieje i wynosi $L$, wybieramy $\delta$ z definicji dla dokładności $\varepsilon/2$. Nierówność trójkąta daje
$d_Y(f(x),f(z))\le d_Y(f(x),L)+d_Y(L,f(z))<\varepsilon$.

Odwrotnie, wybierzmy z lematu \ref{ma1:lem:accumulation-sequence} ciąg $(x_n)$ w $A\setminus\{p\}$ zbieżny do $p$. Dla każdej $\varepsilon>0$ istnieje $\delta$ z założenia, a od pewnego indeksu wszystkie $x_n$ mają odległość od $p$ mniejszą niż $\delta$. Stąd $(f(x_n))$ jest ciągiem Cauchy'ego. Zupełność $Y$ daje punkt $L\in Y$ taki, że $f(x_n)\to L$.

Ustalmy $\varepsilon>0$ i wybierzmy $\delta$ z założenia dla dokładności $\varepsilon/2$. Istnieje indeks $n_0$ taki, że $0<d_X(x_{n_0},p)<\delta$ i $d_Y(f(x_{n_0}),L)<\varepsilon/2$. Dla dowolnego $x\in A$ z $0<d_X(x,p)<\delta$ otrzymujemy
$d_Y(f(x),L)\le d_Y(f(x),f(x_{n_0}))+d_Y(f(x_{n_0}),L)<\varepsilon$.
To jest definicja granicy.
\end{proof}
\begin{przyklad}
\label{ma1:ex:incomplete-limit-target}
Zupełność przestrzeni wartości jest istotna. Niech $X=\R$, $A=(0,1)$, $Y=(0,1)$ z metryką odziedziczoną z $\R$, a $f(x)=x$. W punkcie $p=0$ warunek Cauchy'ego zachodzi: jeśli $0<x,z<\delta$, to $|x-z|<\delta$. Granica w $Y$ nie istnieje, ponieważ ciąg $x_n=1/(n+1)$ ma w $\R$ granicę $0\notin Y$. Gdyby granica w $Y$ istniała, ten sam ciąg miałby w $\R$ dwie różne granice.
\end{przyklad}

\section{Rachunek granic i nierówności}
\begin{twierdzenie}
\label{ma1:thm:function-limit-algebra}
Niech $(X,d_X)$ będzie przestrzenią metryczną, $A\subseteq X$, $p\in A'$, $K=\R$ lub $K=\C$ oraz $f,g:A\to K$. Jeżeli $\lim_{x\to p}f(x)=a$ i $\lim_{x\to p}g(x)=b$, to dla dowolnych $\alpha,\beta\in K$
\[
 \lim_{x\to p}(\alpha f+\beta g)(x)=\alpha a+\beta b,
 \qquad \lim_{x\to p}(fg)(x)=ab.
\]
Ponadto $\lim_{x\to p}|f(x)|=|a|$. Jeżeli $b\ne0$, to istnieje $\delta>0$ takie, że $g(x)\ne0$ dla $x\in A$ z $0<d_X(x,p)<\delta$, a iloraz $f/g$ określony na $B=\{x\in A:g(x)\ne0\}$ ma granicę $a/b$ w $p$ względem $B$.
\end{twierdzenie}
\begin{proof}
Dla dowolnego ciągu $x_n\in A\setminus\{p\}$ zbieżnego do $p$ mamy $f(x_n)\to a$ oraz $g(x_n)\to b$. Rachunek granic ciągów w $K$ daje granicę kombinacji liniowej, iloczynu i modułu. Z kryterium \ref{ma1:thm:function-limit-sequential} otrzymujemy pierwsze trzy tezy.

Jeżeli $b\ne0$, to z definicji granicy dla dokładności $|b|/2$ otrzymujemy $\delta>0$ takie, że $|g(x)-b|<|b|/2$ w przekłutym otoczeniu. Z nierówności odwrotnej $|g(x)|\ge|b|-|g(x)-b|>|b|/2$. W szczególności $B$ zawiera wszystkie punkty $A$ dostatecznie bliskie $p$, różne od $p$, więc $p\in B'$. Dla każdego ciągu w $B\setminus\{p\}$ zbieżnego do $p$ granica ilorazu obrazów wynosi $a/b$ z rachunku granic ciągów. Ponownie stosujemy kryterium ciągowe.
\end{proof}

\begin{twierdzenie}
\label{ma1:thm:function-limit-order}
Niech $(X,d_X)$ będzie przestrzenią metryczną, $A\subseteq X$, $p\in A'$, $f,g,h:A\to\R$. Załóżmy, że istnieje $\eta>0$ takie, że dla $x\in A$ z $0<d_X(x,p)<\eta$ zachodzi $f(x)\le g(x)\le h(x)$.
\begin{enumerate}[label=(\roman*)]
\item Jeżeli istnieją skończone granice $a=\lim_{x\to p}f(x)$ i $b=\lim_{x\to p}g(x)$, to $a\le b$.
\item Jeżeli $\lim_{x\to p}f(x)=\lim_{x\to p}h(x)=L\in\R$, to $\lim_{x\to p}g(x)=L$.
\end{enumerate}
\end{twierdzenie}
\begin{proof}
Dla każdego ciągu $x_n\in A\setminus\{p\}$ dążącego do $p$ założone nierówności zachodzą od pewnego indeksu. W (i) wybieramy jeden taki ciąg, którego istnienie wynika z lematu \ref{ma1:lem:accumulation-sequence}, i stosujemy zachowanie nieostrej nierówności w granicy ciągów rzeczywistych. W (ii) twierdzenie o trzech ciągach daje $g(x_n)\to L$ dla każdego takiego ciągu. Kryterium ciągowe daje tezę.
\end{proof}

\begin{przyklad}
\label{ma1:ex:limit-polynomial-rational}
Dla $p\in\R$ mamy $\lim_{x\to p}x=p$ oraz granicę funkcji stałej równą jej wartości. Pierwszy fakt wynika z wyboru $\delta=\varepsilon$, drugi z zerowej różnicy od granicy. Rachunek granic daje przez indukcję $\lim_{x\to p}x^k=p^k$ dla całkowitych $k\ge0$ i $\lim_{x\to p}P(x)=P(p)$ dla wielomianu $P$. Dla $x\ne2$ mamy $(x^2-4)/(x-2)=x+2$, zatem $\lim_{x\to2}(x^2-4)/(x-2)=4$, mimo że pierwotny iloraz nie jest określony w $2$.
\end{przyklad}
\begin{przyklad}
\label{ma1:ex:limit-squeeze}
Jeżeli $u:\R\setminus\{0\}\to\R$ spełnia $|u(x)|\le1$ dla każdego $x\ne0$, to $-|x|\le xu(x)\le |x|$. Ponieważ $|x|\to0$ przy $x\to0$, twierdzenie \ref{ma1:thm:function-limit-order} daje $\lim_{x\to0}xu(x)=0$, niezależnie od tego, czy $u$ ma granicę w zerze.
\end{przyklad}

\section{Granice jednostronne}
\begin{definicja}
\label{ma1:def:one-sided-limits}
Niech $A\subseteq\R$, $f:A\to Y$, gdzie $Y$ jest przestrzenią metryczną, oraz $p\in\R$. Niech $A_-=A\cap(-\infty,p)$ i $A_+=A\cap(p,+\infty)$. Jeżeli $p\in A_-'$, to granicą lewostronną $\lim_{x\to p^-}f(x)$ nazywamy granicę względem $A_-$. Jeżeli $p\in A_+'$, to granicą prawostronną $\lim_{x\to p^+}f(x)$ nazywamy granicę względem $A_+$.
\end{definicja}
\begin{twierdzenie}
\label{ma1:thm:one-sided-limit-equivalence}
Niech $A\subseteq\R$, $f:A\to Y$, gdzie $Y$ jest przestrzenią metryczną, $p\in (A\cap(-\infty,p))'\cap(A\cap(p,+\infty))'$ oraz $L\in Y$. Wtedy $\lim_{A\ni x\to p}f(x)=L$ wtedy i tylko wtedy, gdy obie granice jednostronne istnieją i są równe $L$.
\end{twierdzenie}
\begin{proof}
Punkt $p$ nie ma wpływu na definicję granicy, więc możemy zastąpić $A$ przez $A\setminus\{p\}=A_-\cup A_+$. Teza wynika z twierdzenia \ref{ma1:thm:limit-restriction-union}.
\end{proof}
\begin{przyklad}
\label{ma1:ex:sign-one-sided}
Niech $f(x)=-1$ dla $x<0$, $f(0)=0$ i $f(x)=1$ dla $x>0$. Granica lewostronna w zerze wynosi $-1$, prawostronna $1$, zatem granica obustronna nie istnieje.
\end{przyklad}
\begin{center}
\begin{tikzpicture}[x=1.5cm,y=.85cm,>=Stealth]
 \draw[->] (-2.1,0)--(2.1,0) node[right] {$x$};
 \draw[->] (0,-1.6)--(0,1.65) node[above] {$y$};
 \draw[thick,blue!65!black] (-1.9,-1)--(0,-1);
 \draw[thick,blue!65!black] (0,1)--(1.9,1);
 \draw[blue!65!black,thick,fill=white] (0,-1) circle (2.3pt);
 \draw[blue!65!black,thick,fill=white] (0,1) circle (2.3pt);
 \fill (0,0) circle (2pt);
 \node[below left] at (0,-1) {$-1$};
 \node[above left] at (0,1) {$1$};
 \node[below right] at (0,0) {$0$};
\end{tikzpicture}

\small Rysunek: różne granice jednostronne funkcji znaku; puste kółka oznaczają wartości nieprzyjmowane w punkcie $0$.
\end{center}

\section{Granice niewłaściwe i granice w nieskończoności}
\begin{definicja}
\label{ma1:def:extended-function-limits}
Niech $A\subseteq\R$, $f:A\to\R$ oraz $p\in A'$. Piszemy $\lim_{A\ni x\to p}f(x)=+\infty$, jeżeli dla każdego $M>0$ istnieje $\delta>0$ takie, że $f(x)>M$ dla wszystkich $x\in A$ z $0<|x-p|<\delta$. Granicę $-\infty$ definiujemy analogicznie, wymagając $f(x)<-M$.

Jeżeli $A$ jest nieograniczony z góry, to $\lim_{A\ni x\to+\infty}f(x)=L\in\R$ oznacza, że dla każdego $\varepsilon>0$ istnieje $R>0$ takie, że $|f(x)-L|<\varepsilon$ dla wszystkich $x\in A$ z $x>R$. Dla granicy przy $x\to-\infty$ wymagamy, aby $A$ był nieograniczony z dołu, i zastępujemy $x>R$ przez $x<-R$. Dla granic niewłaściwych przy $x\to\pm\infty$ zastępujemy warunek $|f(x)-L|<\varepsilon$ odpowiednio przez $f(x)>M$ lub $f(x)<-M$, dla każdego $M>0$.
\end{definicja}
Granice niewłaściwe jednostronne definiujemy przez ograniczenie do odpowiedniego zbioru z definicji \ref{ma1:def:one-sided-limits}. W granicach w nieskończoności założenie o nieograniczoności dziedziny pełni rolę warunku punktu skupienia: dla każdego progu istnieją argumenty leżące poza nim.

\begin{twierdzenie}
\label{ma1:thm:extended-limits-sequential}
Niech $A\subseteq\R$, $f:A\to\R$. Niech $p\in\R\cup\{-\infty,+\infty\}$ będzie taki, że dla skończonego $p$ mamy $p\in A'$, a dla $p=+\infty$ (odpowiednio $p=-\infty$) zbiór $A$ jest nieograniczony z góry (odpowiednio z dołu). Dla $L\in\R\cup\{-\infty,+\infty\}$ granica $\lim_{A\ni x\to p}f(x)=L$ zachodzi wtedy i tylko wtedy, gdy dla każdego ciągu $(x_n)$ w $A$ dążącego do $p$, przy dodatkowym warunku $x_n\ne p$ dla skończonego $p$, mamy $f(x_n)\to L$.
\end{twierdzenie}
\begin{proof}
Dla skończonych $p,L$ teza jest kryterium \ref{ma1:thm:function-limit-sequential}. W pozostałych przypadkach rozumowanie ma tę samą postać. Niech
\[
 U_n=\begin{cases}
 \{x\in A:0<|x-p|<1/n\}, &p\in\R,\\
 \{x\in A:x>n\}, &p=+\infty,\\
 \{x\in A:x<-n\}, &p=-\infty.
 \end{cases}
\]
Wszystkie te zbiory są niepuste na mocy założeń. Warunki definicyjne oznaczają, że dla każdego dopuszczalnego zbioru wartości $V$, postaci $(L-\varepsilon,L+\varepsilon)$ dla skończonego $L$, $(M,+\infty)$ dla $L=+\infty$ lub $(-\infty,-M)$ dla $L=-\infty$, istnieje $n$ takie, że $f(U_n)\subseteq V$. Równoważność wynika z możliwości wyboru $1/n<\delta$ albo $n>R$; w drugą stronę progi $1/n,n$ są dopuszczalne w definicji.

Jeżeli warunek zachodzi, to każdy dopuszczalny ciąg $x_k\to p$ od pewnego indeksu należy do wybranego $U_n$, więc jego obrazy należą do $V$. Odwrotnie, jeśli warunek nie zachodzi, to istnieje jeden zbiór wartości $V$, dla którego dla każdego $n$ wybieramy $x_n\in U_n$ z $f(x_n)\notin V$. Z konstrukcji $x_n\to p$, lecz obrazy nie dążą do $L$, co przeczy założeniu.
\end{proof}
\begin{przyklad}
\label{ma1:ex:extended-limits}
Dla $f(x)=1/x$ na $\R\setminus\{0\}$ i $M>0$ wybieramy $\delta=1/M$. Jeśli $0<x<\delta$, to $1/x>M$; jeśli $-\delta<x<0$, to $1/x<-M$. Zatem $\lim_{x\to0^+}1/x=+\infty$ i $\lim_{x\to0^-}1/x=-\infty$. Granica obustronna nie istnieje nawet w rozszerzonym zbiorze liczb rzeczywistych, co wynika z kryterium ciągowego dla ciągów $1/n$ i $-1/n$.

Natomiast dla $\varepsilon>0$ wybierając $R=1/\varepsilon$ otrzymujemy $|1/x|<\varepsilon$ dla $x>R$ oraz dla $x<-R$. Stąd obie granice w nieskończoności wynoszą $0$. Dla $g(x)=1/x^2$ wybór $\delta=1/\sqrt M$ daje $g(x)>M$ dla $0<|x|<\delta$, czyli granicę obustronną $+\infty$ w zerze.
\end{przyklad}

W następnym rozdziale przejdziemy od zachowania odwzorowania w przekłutym otoczeniu punktu do porównania jego granicy z wartością w tym punkcie, czyli do ciągłości.
